\documentclass[UTF8,zihao=-4]{ctexart}

\usepackage[a4paper,margin=2.4cm]{geometry}
\usepackage{array}
\usepackage{booktabs}
\usepackage{caption}
\usepackage{enumitem}
\usepackage{float}
\usepackage{graphicx}
\usepackage{hyperref}
\usepackage{longtable}
\usepackage{tabularx}
\usepackage{xcolor}

\graphicspath{
  {项目管理作业2/项目管理作业2/城东酥肉股份有限公司-项目管理作业1/}
  {城东酥肉股份有限公司-项目管理作业1/}
  {项目管理作业2/项目管理作业2/}
}

\hypersetup{
  colorlinks=true,
  linkcolor=black,
  urlcolor=blue,
  pdftitle={城东酥肉股份有限公司韩餐店两阶段项目管理汇报书},
  pdfauthor={项目组}
}

\setlist[itemize]{leftmargin=2em,itemsep=0.2em,topsep=0.2em}
\setlist[enumerate]{leftmargin=2em,itemsep=0.2em,topsep=0.2em}
\renewcommand{\arraystretch}{1.25}

\newcolumntype{Y}{>{\raggedright\arraybackslash}X}
\newcommand{\maybeimage}[3]{
  \IfFileExists{#1}{
    \begin{figure}[H]
      \centering
      \includegraphics[width=#2\textwidth]{#1}
      \caption{#3}
    \end{figure}
  }{}
}

\title{\heiti 城东酥肉股份有限公司“首尔之春”韩餐店\\两阶段建设项目管理汇报书}
\author{项目组}
\date{\today}

\begin{document}

\maketitle

\begin{abstract}
本汇报书围绕“首尔之春”韩餐店建设项目展开，采用“两阶段投资、先验证后优化”的项目管理思路：项目全程以200万元为资金上限，前期建设小型韩餐店，验证选址、菜品、供应链、人员培训和校园客群转化；后期在运营指标达标后，继续围绕200万元预算进行门店优化、品牌露出和标准化运营能力建设。报告从项目范围、进度计划、关键路径、PERT三点估算、挣值分析、资源重排、装修外包、风险控制和阶段门决策等方面进行说明，目标是在控制成本与工期的前提下完成门店开业，并为后续扩张提供可复用的管理模板。
\end{abstract}

\tableofcontents
\newpage

\section{项目背景与总体思路}

\subsection{项目背景}

随着大学城周边餐饮消费持续增长，韩式烤肉、韩式拌饭、部队锅、炸鸡年糕等产品兼具社交属性和高频消费潜力，适合面向学生、教职工和周边居民开展中档餐饮服务。本项目拟以“首尔之春”为门店品牌，依托城东酥肉股份有限公司现有组织与资金基础，完成韩餐店从策划、选址、证照、装修、设备、供应链、招聘培训到试营业和正式开业的全过程管理。

项目采用“两阶段”路线：第一阶段投入200万元建设小型韩餐店，降低一次性投入风险，集中验证产品、客流和运营流程；第二阶段在第一阶段经营数据达标后，继续以200万元预算为上限，优化门店配置、品牌曝光和供应链议价能力。

\subsection{两阶段建设目标}

\begin{table}[H]
  \centering
  \caption{两阶段项目目标对比}
  \begin{tabularx}{\textwidth}{>{\bfseries}p{2.4cm}YY}
    \toprule
    项目维度 & 前期：小型韩餐店 & 后期：优化型韩餐店 \\
    \midrule
    投资规模 & 200万元 & 200万元 \\
    战略定位 & 低风险试点店，验证商业模式 & 优化提升店，强化品牌与运营能力 \\
    目标客群 & 大学生、教职工、周边居民 & 校园客群、周边居民、家庭聚餐和小型团餐客流 \\
    管理重点 & 控制预算、快速开业、形成标准作业流程 & 优化产能、提升品牌、完善供应链和人才梯队 \\
    核心交付物 & 小型门店开业、基础菜单、供应商合同、培训记录、试营业复盘 & 优化后门店开业、品牌视觉、采购清单、管理制度、营销体系 \\
    成功标准 & 门店稳定运营，核心菜品毛利达标，试营业反馈良好 & 客流和翻台率稳定，品牌知名度提升，具备复制条件 \\
    \bottomrule
  \end{tabularx}
\end{table}

\section{项目范围与交付成果}

\subsection{项目范围说明}

本项目范围包括韩餐店开业前后必要的全部项目活动，主要覆盖以下五类工作包：

\begin{enumerate}
  \item 启动与策划：立项分工、市场调研、商业计划、预算测算、选址谈判与租赁签约。
  \item 证照与装修：装修设计、营业执照、食品经营许可、装修施工、墙面粉刷外包、消防与卫生验收。
  \item 菜品与供应链：菜单研发、价格测算、设备采购、设备安装调试、食材供应商签约与开业备货。
  \item 招聘与培训：店长、厨师、前厅、后厨和兼职人员招聘，服务流程培训、岗前试菜、排班和应急预案。
  \item 推广与开业：校园社群预热、试营业、正式开业活动、首月经营数据复盘。
\end{enumerate}

\subsection{不纳入本期范围的内容}

第一阶段小型店不建设中央厨房，不进行跨城市连锁扩张，不采购明显超出门店产能的重资产设备。第二阶段大型店是否建设中央厨房、会员系统和跨店配送能力，应在第一阶段运营数据达标后另行立项。

\subsection{主要交付成果}

\begin{table}[H]
  \centering
  \caption{项目主要交付成果}
  \begin{tabularx}{\textwidth}{p{3cm}Y}
    \toprule
    交付类别 & 具体成果 \\
    \midrule
    管理类成果 & 项目章程、WBS、进度计划、责任矩阵、风险清单、阶段复盘报告 \\
    经营类成果 & 商业计划书、菜单与定价表、供应商合同、采购清单、库存标准 \\
    建设类成果 & 租赁合同、装修设计图、施工验收记录、设备安装调试记录 \\
    人员类成果 & 招聘结果、培训记录、岗位说明书、排班表、服务标准流程 \\
    开业类成果 & 试营业记录、客户反馈、营销物料、正式开业方案、首月复盘报告 \\
    \bottomrule
  \end{tabularx}
\end{table}

\maybeimage{项目管理作业2/项目管理作业2/城东酥肉股份有限公司-项目管理作业1/WBS分解图.png}{0.86}{项目WBS分解图}

\section{组织结构与责任分工}

\subsection{项目组织}

项目组设项目负责人1名，统筹范围、进度、预算与风险；成员A负责商业计划、预算和选址谈判；成员B负责装修设计、证照跟进和装修验收；成员C负责设备采购、安装调试、供应商签约和备货；成员D负责招聘培训、营销预热、试营业和客户反馈。

\subsection{RACI责任矩阵}

\begin{table}[H]
  \centering
  \caption{核心工作RACI责任矩阵}
  \begin{tabularx}{\textwidth}{p{4cm}ccccY}
    \toprule
    工作内容 & A & B & C & D & 说明 \\
    \midrule
    商业计划与预算 & A/R & C & C & I & A负责测算投资回报和预算边界 \\
    选址谈判签约 & A/R & C & I & I & 与租金、客流和开业窗口直接相关 \\
    装修设计确认 & C & A/R & C & I & B对方案和施工接口最终负责 \\
    墙面粉刷外包 & I & A & C & I & 外包队为R，B负责验收闭环 \\
    设备采购与安装 & I & C & A/R & I & C负责设备供应商与现场调试 \\
    招聘与培训 & I & I & C & A/R & D负责岗位到位和服务流程培训 \\
    试营业与开业活动 & C & I & C & A/R & D组织试营业反馈和开业营销 \\
    \bottomrule
  \end{tabularx}
\end{table}

注：A表示最终负责，R表示执行，C表示协作或咨询，I表示知会。

\maybeimage{项目管理作业2/项目管理作业2/城东酥肉股份有限公司-项目管理作业1/责任矩阵图.png}{0.78}{项目责任矩阵图}

\section{进度计划与关键路径}

\subsection{基准工期}

根据现有MPP计划，项目基准总工期为112天。项目从立项分工开始，经过市场调研、商业计划、选址签约、装修设计、装修施工、设备安装、供应商签约、试营业和正式开业，最终完成门店开业交付。

\subsection{关键路径}

项目关键路径为：

\begin{center}
  \fbox{\parbox{0.92\textwidth}{
  A1立项分工 $\rightarrow$ A2市场调研 $\rightarrow$ A3商业计划与预算 $\rightarrow$ A4选址谈判签约 $\rightarrow$ B1装修设计 $\rightarrow$ B3装修施工 $\rightarrow$ C3设备安装调试 $\rightarrow$ E1/E2供应商签约与备货 $\rightarrow$ F2试营业 $\rightarrow$ H开业前缓冲 $\rightarrow$ G正式开业阶段
  }}
\end{center}

其中，B3装修施工工期为20天，是连接选址、设计、设备和试营业的关键环节；H开业前缓冲为5天，用于试营业反馈、现场收口、物料确认和突发问题处理，不宜随意取消。

\subsection{关键活动控制要点}

\begin{longtable}{p{2cm}p{4cm}p{8cm}}
  \caption{关键活动控制要点}\\
  \toprule
  编号 & 活动 & 管理重点 \\
  \midrule
  A4 & 选址谈判签约 & 控制租金水平、免租期、物业条件、排烟和消防可行性，避免后期返工。 \\
  B1 & 装修设计 & 在第35至45天完成方案确认，保障第45天装修施工进场。 \\
  B3 & 装修施工 & 严控水电、排烟、消防、墙面粉刷、通风保护和收口验收。 \\
  C3 & 设备安装调试 & 设备进场前确认电力、排水、排烟和操作动线。 \\
  E1/E2 & 供应签约与备货 & 先锁定供应框架、报价和账期，再组织开业备货。 \\
  F2 & 试营业 & 检验出餐效率、服务流程、菜品稳定性和顾客反馈。 \\
  G & 正式开业 & 控制开业活动强度，避免营销承诺超过门店承载能力。 \\
  \bottomrule
\end{longtable}

\maybeimage{项目管理作业2/项目管理作业2/城东酥肉股份有限公司-项目管理作业1/甘特图预览.png}{0.95}{项目甘特图预览}
\maybeimage{项目管理作业2/项目管理作业2/城东酥肉股份有限公司-项目管理作业1/网络图预览.png}{0.95}{项目网络图预览}

\section{PERT三点估算与工期不确定性}

\subsection{估算方法}

为考虑选址、装修、证照和开业活动的不确定性，项目采用PERT三点估算方法。单项活动期望时间计算公式为：

\[
T_e=\frac{a+4m+b}{6}
\]

其中，$a$为最乐观工期，$m$为最可能工期，$b$为最悲观工期。活动方差为：

\[
\sigma^2=\left(\frac{b-a}{6}\right)^2
\]

\subsection{估算结果}

在关键路径活动上进行PERT估算后，项目期望工期为115.33天，约为115.3天；项目方差为32.0，标准差为5.66天。与MPP确定性计划112天相比，考虑风险后的期望工期增加约3.3天。

\begin{table}[H]
  \centering
  \caption{项目完成概率区间}
  \begin{tabular}{ccc}
    \toprule
    置信区间 & 工期范围 & 管理含义 \\
    \midrule
    68\% & 109.7至121.0天 & 常规课堂汇报和内部进度沟通可采用 \\
    95\% & 104.0至126.6天 & 对外承诺开业日期时应重点参考 \\
    99\% & 98.4至132.3天 & 用于极端风险下的压力测试 \\
    \bottomrule
  \end{tabular}
\end{table}

估算结果说明，项目表面上可以按112天推进，但实际管理中应以115天左右作为更稳妥的期望工期，并对选址谈判、装修施工和正式开业活动设置重点监控。

\section{投资预算与成本管理}

\subsection{前期小型店200万元预算}

第一阶段预算以“能开业、能运营、能验证”为原则，避免过度装修和重资产投入。

\begin{table}[H]
  \centering
  \caption{小型韩餐店200万元预算分配}
  \begin{tabularx}{\textwidth}{p{4cm}rY}
    \toprule
    费用项 & 金额（万元） & 用途说明 \\
    \midrule
    租金、押金及转让相关费用 & 35 & 覆盖门店租赁启动成本和必要押金 \\
    装修与基础施工 & 55 & 水电、排烟、墙面、地面、消防和基础装饰 \\
    厨房与前厅设备 & 45 & 烤炉、冷柜、排烟设备、收银系统和餐具 \\
    首批食材与供应链启动 & 15 & 试营业和开业初期备货 \\
    证照、设计与咨询 & 8 & 营业执照、食品经营许可、设计咨询等 \\
    招聘、培训与开办人工 & 18 & 核心岗位招聘、培训和试营业人工 \\
    开业营销 & 12 & 校园社群、团购平台、开业物料和优惠活动 \\
    预备费 & 12 & 应对装修变更、设备补采和证照延迟 \\
    \midrule
    合计 & 200 &  \\
    \bottomrule
  \end{tabularx}
\end{table}

\subsection{后期门店200万元预算}

第二阶段预算仍以200万元为资金上限，在第一阶段数据达标后启动，重点放在门店优化、品牌露出和标准化能力建设，而不是一次性扩大投资规模。

\begin{table}[H]
  \centering
  \caption{后期韩餐店200万元预算分配}
  \begin{tabularx}{\textwidth}{p{4cm}rY}
    \toprule
    费用项 & 金额（万元） & 用途说明 \\
    \midrule
    场地租赁、押金及商圈进入成本 & 36 & 覆盖门店租赁、押金和必要位置优化成本 \\
    设计装修与工程施工 & 60 & 控制装修标准，优先保障排烟消防、墙面地面和工程验收 \\
    厨房设备与前厅设备 & 36 & 补充烤肉、冷柜、洗消、排烟、叫号和收银相关设备 \\
    品牌视觉、IT系统与会员工具 & 10 & 门头、菜单、会员工具、数据记录和线上运营工具 \\
    供应链建设与首批库存 & 16 & 核心食材协议、开业备货和安全库存 \\
    人员招聘与培训 & 16 & 店长、厨师、前厅服务和兼职人员的招聘培训 \\
    开业营销与品牌推广 & 14 & 校园推广、团购平台、探店内容和会员拉新 \\
    预备费 & 12 & 应对工程变更、设备涨价、审批延迟和营销调整 \\
    \midrule
    合计 & 200 &  \\
    \bottomrule
  \end{tabularx}
\end{table}

\subsection{成本控制机制}

项目成本控制采用“预算基线+变更审批+挣值分析”的方式。装修、设备和营销是最容易超支的三类费用，应设置审批阈值：单项变更超过预算10\%或金额超过5万元时，必须由项目负责人复核；涉及关键路径的变更还需同步评估对工期的影响。

\subsection{200万元口径财务测算}

由于本项目资金上限为200万元，财务测算应采用小型门店口径，不能沿用大型店的座位数、客单价和活动规模。按60座、1.8次翻台、68元客单价、330个营业日估算，第一年堂食收入约242.4万元；外卖按60单/天、45元/单、330天估算，约89.1万元；饮品和小食加购约10.7万元，团购及校园活动约38.4万元。因此，第一年营业收入合计约380.6万元。

成本端同步按小型店规模测算。第一年食材成本按收入的38\%估算，人工约78万元，租金约42万元，水电能耗约18万元，营销费用按收入的7.5\%估算，管理及杂费约18万元，税费按收入的5\%估算。由此得到第一年总成本约348.2万元，经营现金流约32.4万元。

\begin{table}[H]
  \centering
  \caption{200万元口径下的财务测算结果}
  \begin{tabular}{lr}
    \toprule
    指标 & 测算结果 \\
    \midrule
    初始投资 & 200.0万元 \\
    第一年营业收入 & 380.6万元 \\
    第一年总成本 & 348.2万元 \\
    第一年经营现金流 & 32.4万元 \\
    NPV（贴现率5\%） & 26.7万元 \\
    静态回本期 & 约4.60年 \\
    贴现后BEP & 第6年 \\
    \bottomrule
  \end{tabular}
\end{table}

测算结果说明，小型店在200万元投资约束下仍具备正NPV，但回本并非第一年完成，而是需要依靠后续经营增长逐步覆盖初始投入。因此，项目管理重点应放在控制装修和设备超支、稳定食材成本率、提高翻台率和复购率，而不是用大型店收入假设放大投资回报。

\section{EVM挣值分析}

\subsection{分析假设}

项目运行到三分之一节点时，计划完成比例为33.3\%，实际进度提前20\%，当前成本超支20\%。据此计算计划价值PV、挣值EV、实际成本AC、成本偏差CV、进度偏差SV、成本绩效指数CPI和进度绩效指数SPI。

\subsection{项目200万元情景}

\begin{table}[H]
  \centering
  \caption{200万元BAC情景下的EVM指标}
  \begin{tabular}{lrl}
    \toprule
    指标 & 数值 & 含义 \\
    \midrule
    BAC & 200.0万元 & 完工预算 \\
    PV & 66.6万元 & 按计划应完成的价值 \\
    EV & 79.9万元 & 实际已完成工作的预算价值 \\
    AC & 95.9万元 & 实际发生成本 \\
    CV & -16.0万元 & 成本超支 \\
    SV & 13.3万元 & 进度提前 \\
    CPI & 0.833 & 成本效率偏低 \\
    SPI & 1.200 & 进度效率高于计划 \\
    \bottomrule
  \end{tabular}
\end{table}

\subsection{EVM管理结论}

EVM结论为：项目进度提前，但成本效率偏低。管理上不应只强调“进度快”，还应关注提前采购、赶工、装修加项和营销预支是否造成预算失控。建议把进度优势转化为风险缓冲，同时冻结非必要变更，优先复核装修、设备和营销支出。

\section{资源受限与进度重排}

\subsection{资源冲突识别}

根据计划安排，资源冲突主要出现在两个区间：

\begin{itemize}
  \item B侧冲突：第35至45天，B1装修设计与B2营业执照、食品经营许可并行，均需要成员B处理设计确认、证照材料和审批沟通。
  \item C侧冲突：第70至75天，E1食材供应商签约与E2开业备货并行，均占用成员C。
\end{itemize}

\subsection{资源重排方案}

资源重排遵循“优先保障关键路径，不轻易拉长总工期”的原则。B侧保持B1第35至45天完成，确保B3第45天进场；B2利用浮时顺延至第45至55天。C侧保持E1第70至75天完成，E2顺延至第75至80天，F2试营业顺延至第80至85天，占用原本的5天开业前缓冲，正式开业仍从第85天开始，总工期仍为112天。

\begin{table}[H]
  \centering
  \caption{资源重排前后对比}
  \begin{tabularx}{\textwidth}{p{3cm}p{3cm}p{3cm}Y}
    \toprule
    活动 & 原计划 & 重排后 & 说明 \\
    \midrule
    B1装修设计 & 第35至45天 & 第35至45天 & 保持不变，保障装修施工进场 \\
    B2证照许可 & 第35至45天 & 第45至55天 & 利用浮时解决B侧冲突 \\
    E1供应商签约 & 第70至75天 & 第70至75天 & 先锁定供应框架和报价 \\
    E2开业备货 & 第70至75天 & 第75至80天 & 顺接供应商签约结果 \\
    F2试营业 & 第75至80天 & 第80至85天 & 占用开业前缓冲 \\
    G正式开业 & 第85至112天 & 第85至112天 & 正式开业节点不变 \\
    \bottomrule
  \end{tabularx}
\end{table}

\maybeimage{项目管理作业2/项目管理作业2/城东酥肉股份有限公司-项目管理作业1/resource-overallocation-chart.png}{0.9}{资源过度分配图}
\maybeimage{项目管理作业2/项目管理作业2/1058.png}{0.9}{资源约束后的进度安排}

\section{装修外包与约束优化}

\subsection{外包决策}

墙面粉刷作为B3装修施工中的专业子项，不由A、B、C、D成员直接执行，而是外包给专业施工队。外包队负责具体施工，成员B负责开工确认、过程抽检、完工验收和整改关闭。

该安排的理由包括：第一，成员B已经参与选址、装修设计、证照和验收，若继续承担粉刷执行，容易形成单点过载；第二，专业粉刷队在基层处理、环保材料、色差控制和工期压缩方面更有经验；第三，将“施工执行”和“验收责任”分开，更利于质量控制。

\subsection{B3装修施工二级WBS}

\begin{table}[H]
  \centering
  \caption{B3装修施工二级WBS}
  \begin{tabularx}{\textwidth}{p{2cm}p{3cm}p{3cm}Y}
    \toprule
    编号 & 子任务 & 时间安排 & 交付要求 \\
    \midrule
    B3.1 & 基础施工 & 第45至51天 & 水电、排烟、消防预留和关键管线完成 \\
    B3.2 & 墙面基层 & 第52至54天 & 找平、修补、打磨，达到粉刷条件 \\
    B3.3 & 墙面粉刷 & 第55至59天 & 外包队执行，B控制质量、进度和材料 \\
    B3.4 & 通风保护 & 第60至61天 & 保证干燥、通风、环保和设备进场条件 \\
    B3.5 & 收口验收 & 第62至65天 & 完成整改清单、验收记录并移交设备安装 \\
    \bottomrule
  \end{tabularx}
\end{table}

\subsection{外包约束}

外包开始时间不是签合同当天，而是设计、预算、材料和现场条件均确认后，具备进场施工条件的日期。外包结束时间也不是工人离场当天，而是施工完成、整改关闭、B验收通过并能移交下一道工序的日期。

为降低返工风险，项目设置三个验收关口：开工前确认墙面颜色、环保材料、报价上限、质保期和完工日期；施工中由B每天或隔天抽检基层平整度、色差、污染保护和现场安全；完工后若验收不通过，由外包队返工，返工费用和延期责任按合同执行。

\section{风险管理}

\begin{longtable}{p{3cm}p{3cm}p{2cm}p{2cm}p{4cm}}
  \caption{项目主要风险与应对措施}\\
  \toprule
  风险事项 & 影响 & 概率 & 影响程度 & 应对措施 \\
  \midrule
  选址谈判延迟 & 推迟装修设计和施工进场 & 中 & 高 & 准备备选铺位，提前核查排烟、消防和租金条件 \\
  证照办理滞后 & 影响试营业和正式营业合规性 & 中 & 高 & 提前准备材料，专人跟进审批节点，必要时咨询专业机构 \\
  装修施工返工 & 增加工期和成本 & 中 & 高 & 明确验收标准，设置过程抽检和10\%变更审批阈值 \\
  设备到货延期 & 影响安装调试和试营业 & 中 & 中 & 关键设备提前下单，准备替代供应商和临时设备方案 \\
  人员招聘不足 & 影响服务质量和翻台效率 & 中 & 中 & 提前招聘店长和厨师长，建立兼职储备池 \\
  开业客流不及预期 & 影响收入和品牌声量 & 中 & 高 & 分阶段预热，试营业收集反馈，优化团购与校园社群投放 \\
  成本超支 & 挤压现金流和营销预算 & 中 & 高 & 按EVM监控CPI，冻结非必要采购，重大变更复核 \\
  食材价格波动 & 影响毛利率 & 中 & 中 & 签订框架协议，设置安全库存和替代菜品方案 \\
  \bottomrule
\end{longtable}

\section{阶段门与后期扩张条件}

\subsection{第一阶段完成标准}

前期200万元小型店不是简单“开出来”，而是要形成后期200万元优化投入的判断依据。第一阶段应在开业后1至3个月完成复盘，重点观察以下指标：

\begin{itemize}
  \item 日均客流、翻台率和高峰期排队情况达到预期；
  \item 核心菜品毛利率稳定，损耗率可控；
  \item 顾客评价集中问题可被菜单、服务或流程调整解决；
  \item 供应商供货稳定，账期、价格和配送质量可控；
  \item 人员排班、培训、出餐和收银流程可复制；
  \item 现金流未出现持续性紧张，装修和设备投入未明显透支。
\end{itemize}

\subsection{进入第二阶段的决策门}

只有当小型店达到阶段门标准后，才建议启动后期200万元门店优化建设。阶段门决策如下：

\begin{table}[H]
  \centering
  \caption{第二阶段启动决策门}
  \begin{tabularx}{\textwidth}{p{3.5cm}p{4cm}Y}
    \toprule
    决策项 & 通过标准 & 决策含义 \\
    \midrule
    经营验证 & 连续经营数据稳定，核心菜品毛利达标 & 证明商业模式可行 \\
    客群验证 & 校园与周边客群复购明显，评价稳定 & 证明品牌有扩张基础 \\
    供应链验证 & 主供应商稳定，备选供应商可用 & 降低后期运营断供风险 \\
    人才验证 & 店长、厨师长和主管岗位可承担培训任务 & 具备复制组织能力 \\
    财务验证 & 现金流和成本偏差在可控区间 & 避免盲目扩大投资 \\
    \bottomrule
  \end{tabularx}
\end{table}

若上述指标未达标，应优先优化第一阶段门店，而不是直接进入后期追加投入。后期200万元项目的本质不是“把店做大”，而是在资金约束内把已经验证的产品、流程和品牌能力做稳。

\section{结论}

本项目以项目管理方法重新组织韩餐店开业工作，将看似零散的选址、装修、证照、采购、招聘、培训和营销活动整合为可计划、可监控、可复盘的项目体系。前期200万元小型店承担商业模式验证任务，重点是控制预算、压缩试错成本、形成标准流程；后期200万元门店优化承担运营提升任务，重点是品牌建设、产能优化、供应链稳定和组织复制。

从进度角度看，MPP基准工期为112天，PERT期望工期为115.3天，说明项目应保留合理缓冲。从成本角度看，EVM显示项目可能出现“进度提前但成本超支”的状态，需要严格控制装修、设备和营销费用。从资源角度看，通过B侧和C侧错峰安排，可以在不增加总工期的前提下消除资源冲突。从外包角度看，墙面粉刷交由专业队伍执行，项目组保留验收和接口控制，可以降低返工与人员过载风险。

因此，本项目建议按照“200万元试点、数据复盘、阶段门评审、200万元优化”的路径推进。这样既能保障前期开业落地，也能为后期韩餐店优化建设提供更可靠的决策依据。

\end{document}
